\subsection{Homological dimension of Noetherian rings}

Let A be a Noetherian local ring, and M a finitely generated A-module. Recall (A, X, § 3, no. 6) that a resolution

\[
\dots \longrightarrow \mathrm{L} _ {n} \stackrel {d _ {n}} {\longrightarrow} \mathrm{L} _ {n - 1} \longrightarrow \dots \longrightarrow \mathrm{L} _ {0} \stackrel {d _ {0}} {\longrightarrow} \mathrm{M} \rightarrow 0
\]

of M is a minimal projective resolution if each of the modules \(L_{i}\) is free of finite type, and if the complex \(\kappa_{A} \otimes_{A} L\) has zero differential. For every integer \(i \geqslant 0\) ,

we have

\[
[ \mathrm{Ext} _ {\mathrm{A}} ^ {i} (\mathrm{M}, \kappa_ {\mathrm{A}}): \kappa_ {\mathrm{A}} ] = [ \mathrm{Tor} _ {i} ^ {\mathrm{A}} (\mathrm{M}, \kappa_ {\mathrm{A}}): \kappa_ {\Lambda} ] = \mathrm{rg} _ {\mathrm{A}} (\mathrm{L} _ {i})\tag{1}
\]

(A, X, p.~103, example 3). Every finitely generated A-module admits such a resolution (A, X, p.~56, prop. 10).

PROPOSITION 4. Let A be a Noetherian local ring, M a finitely generated A-module, and n an integer ≥ 0. The following conditions are equivalent:

\begin{enumerate}
\def\labelenumi{(\roman{enumi})}
\item
  we have \(\mathrm{dp}_{\mathrm{A}}(\mathrm{M}) < n\)
\item
  we have \(\mathrm{Tor}_n^{\mathrm{A}}(\mathrm{M},\kappa_{\mathrm{A}}) = 0\) ;
\item
  we have \(\operatorname{Ext}_{\mathrm{A}}^{n}(\mathrm{M},\kappa_{\mathrm{A}}) = 0\) ;
\item
  every minimal projective resolution of M has length \textless{} n.
\end{enumerate}

The implications (i)⇒(ii) and (i)⇒(iii) are immediate (A, X, p.~100, th. 1). Let L be a minimal projective resolution of M; if (ii) or (iii) holds, then we have \(L_{n}=0\) by (1). Since every minimal projective resolution of M is isomorphic to L (A, X, p.~54, prop. 8), we deduce (iv). The implication (iv)⇒(i) is trivial.

COROLLARY 1.--- Let A be a Noetherian local ring and n an integer ≥ 0. The following conditions are equivalent:

\begin{enumerate}
\def\labelenumi{(\roman{enumi})}
\item
  we have \(\mathrm{dh}(\mathrm{A}) < n\)
\item
  we have \(\mathrm{Ext}_{A}^{i}(\mathrm{M},\mathrm{N}) = 0\) and \(\mathrm{Tor}_i^A (\mathrm{M},\mathrm{N}) = 0\) for every pair \((\mathbf{M},\mathbf{N})\) of A-modules and every integer \(i\geqslant n\) ;
\item
  we have \(\mathrm{Tor}_n^{\mathrm{A}}(\kappa_{\mathrm{A}},\kappa_{\mathrm{A}}) = 0\)
\item
  we have \(\operatorname{Ext}_{\mathrm{A}}^{n}(\kappa_{\mathrm{A}},\kappa_{\mathrm{A}}) = 0\) ;
\item
  we have \(\mathrm{dp}_{\mathrm{A}}(\kappa_{\mathrm{A}}) < n\) .
\end{enumerate}

It is clear that (i) implies (ii) and that (ii) implies (iii) and (iv). By Prop. 4 applied to the A-module \(\kappa_{\mathrm{A}}\) each of the conditions (iii) and (iv) implies (v). We prove that (v) implies (i): if \(\mathrm{dp}_{A}(\kappa_{\mathrm{A}}) < n\) , one has \(\operatorname{Tor}_{n}^{\mathrm{A}}(\mathrm{M}, \kappa_{\mathrm{A}}) = 0\) for every A-module M; consequently every finitely generated A-module has finite projective dimension \(< n\) (prop. 4), which entails \(\mathrm{dh}(\mathrm{A}) < n\) (A, X, p. 138, prop. 4).

\[
\text { COROLLARY   2. } - \text { We   have   dh(A) = dp } _ {A} (\kappa_ {A})  .
\]

Remarks.- 1) Let A be a local ring. The A-module \(\mathrm{Tor}_{1}^{\mathrm{A}}(\kappa_{\mathrm{A}},\kappa_{\mathrm{A}})\) is isomorphic to \(\mathfrak{m}_{A}/\mathfrak{m}_{A}^{2}\). Consequently, when A is Noetherian, for \(\mathrm{Tor}_{1}^{\mathrm{A}}(\kappa_{\mathrm{A}},\kappa_{\mathrm{A}})\) to be zero, it is necessary and sufficient that \(\mathfrak{m}_{\mathrm{A}}\) be zero, that is, that A be a field. Cor. 1 therefore implies that a Noetherian local ring of homological dimension 0 is a field.

\begin{enumerate}
\def\labelenumi{\arabic{enumi})}
\setcounter{enumi}{1}
\tightlist
\item
  Let A be a Noetherian local ring, M a finitely generated A-module of finite projective dimension n, and N a nonzero finitely generated A-module. The A-module Ext \(_{A}^{n}(M,N)\) is nonzero: indeed, let L be a minimal projective resolution of M, and let d be its differential. We have an exact sequence
\end{enumerate}

\[
\operatorname{Hom} _ {\mathrm{A}} \left(\mathrm{L} _ {n - 1}, \mathrm{N}\right) \xrightarrow {\operatorname{Hom} \left(d _ {n} , 1\right)} \operatorname{Hom} _ {\mathrm{A}} \left(\mathrm{L} _ {n}, \mathrm{N}\right) \longrightarrow \operatorname{Ext} _ {\mathrm{A}} ^ {n} (\mathrm{M}, \mathrm{N}) \rightarrow 0.
\]

As \(d_n \otimes 1_{\kappa_A}\) is zero, we deduce by tensor product with \(\kappa_A\) an isomorphism \(\kappa_A \otimes_A \operatorname{Hom}_A(L_n, N) \longrightarrow \kappa_A \otimes_A \operatorname{Ext}_A^n(M, N)\) , whence, taking into account formula (1) above,

\[
\left[ \kappa_ {\mathrm{A}} \otimes_ {\mathrm{A}} \operatorname{Ext} _ {\mathrm{A}} ^ {n} (\mathrm{M}, \mathrm{N}): \kappa_ {\mathrm{A}} \right] = \left[ \operatorname{Ext} _ {\mathrm{A}} ^ {n} (\mathrm{M}, \kappa_ {\mathrm{A}}): \kappa_ {\mathrm{A}} \right] \left[ \kappa_ {\mathrm{A}} \otimes_ {\mathrm{A}} \mathrm{N}: \kappa_ {A} \right],
\]

which is nonzero by Prop. 4 and Nakayama’s lemma. Consequently, the projective dimension of M is the largest integer i such that Ext \(_{A}^{i}(M,N)\) is nonzero.

\begin{enumerate}
\def\labelenumi{\arabic{enumi})}
\setcounter{enumi}{2}
\tightlist
\item
  Let A be a Noetherian ring, M a finitely generated A-module of finite projective dimension, and N a finitely generated A-module whose support is equal to \(\operatorname{Spec}(A)\) . By the preceding remark and Propositions 2 and 3 of No. 2, the projective dimension \(n\) of \(M\) is the largest integer \(i\) such that \(\operatorname{Ext}_{A}^{i}(M,N)\) is nonzero; the support of the A-module \(\operatorname{Ext}_{A}^{n}(M,N)\) is the set of elements \(\mathfrak{p}\) of \(\operatorname{Spec}(A)\) such that \(\mathrm{dp}_{A_{\mathfrak{p}}}(M_{\mathfrak{p}})=n\) .
\end{enumerate}

There may exist finitely generated A-modules M of finite projective dimension \(+\infty\) , satisfying \(\operatorname{Ext}_{\mathrm{A}}^{i}(\mathrm{M},\mathrm{A})=0\) for \(i\) large enough: * this is the case, for example, of the A-module \(\kappa_{\Lambda}\) when \(\Lambda\) is a Gorenstein local ring that is not regular*.

PROPOSITION 5.--- Let A be a Noetherian ring, M a finitely generated A-module, and n an integer ≥ 0. The following conditions are equivalent:

\begin{enumerate}
\def\labelenumi{(\roman{enumi})}
\item
  we have \(\mathrm{dp}_{\mathrm{A}}(\mathrm{M}) < n\)
\item
  for every maximal ideal \(\mathfrak{m}\) of \(A\) , one has \(\operatorname{Ext}_{\mathrm{A}}^{n}(\mathrm{M},\mathrm{A}/\mathfrak{m}) = 0\) (respectively, one has \(\operatorname{Tor}_{n}^{\mathrm{A}}(\mathrm{M},\mathrm{A}/\mathfrak{m}) = 0\) );
\item
  for every maximal ideal m of A, we have \(\operatorname{Ext}_{\mathrm{A}_{m}}^{n}(\mathrm{M}_{\mathfrak{m}},\mathrm{A}/\mathfrak{m})=0\) (respectively, one has \(\operatorname{Tor}_{n}^{\mathrm{A}_{\mathfrak{m}}}(\mathrm{M}_{\mathfrak{m}},A/\mathfrak{m})=0\) ).
\end{enumerate}

(i)⇒(ii): this is clear.

(ii)⇒(iii): this follows from prop. 2 of no. 2.

\begin{enumerate}
\def\labelenumi{(\roman{enumi})}
\setcounter{enumi}{2}
\tightlist
\item
  (iii) \(\Rightarrow\) (i): by prop. 4, condition (iii) implies the inequality \(\mathrm{dp}_{A_{\mathfrak{m}}}(M_{\mathfrak{m}}) < n\) for every maximal ideal \(\mathfrak{m}\) of A, and we conclude by Prop. 3.
\end{enumerate}

Remark 4. — Let A be a Noetherian ring and \(n\) an integer \(\geqslant 0\). If \(\mathfrak{m}\) and \(\mathfrak{m}'\) are two distinct maximal ideals of A, the A-modules \(\mathrm{Ext}_{\mathrm{A}}^{n}(\mathrm{A/m},\mathrm{A/m}')\) and \(\mathrm{Tor}_{n}^{\mathrm{A}}(\mathrm{A/m},\mathrm{A/m}')\) are annihilated by \(\mathfrak{m} + \mathfrak{m}'\) and therefore are zero. By a proof analogous to that of Cor. 1 of Prop. 4, one deduces from Prop. 5 the equivalence of the following conditions:

\begin{enumerate}
\def\labelenumi{(\roman{enumi})}
\item
  we have \(\mathrm{dh}(\mathrm{A}) < n\)
\item
  we have \(\operatorname{Ext}_A^i (\mathrm{M},\mathrm{N}) = 0\) and \(\operatorname{Tor}_i^{\mathrm{A}}(\mathrm{M},\mathrm{N}) = 0\) for every pair \((\mathrm{M},\mathrm{N})\) of A-modules and every integer \(i\geqslant n\) ;
\item
  we have \(\mathrm{Tor}_n^{\mathrm{A}}(\mathrm{A} / \mathfrak{m},\mathrm{A} / \mathfrak{m}) = 0\) for every maximal ideal \(\mathfrak{m}\) of A;
\item
  we have \(\operatorname{Ext}_{\mathsf{A}}^{n}(\mathsf{A} / \mathfrak{m},\mathsf{A} / \mathfrak{m}) = 0\) for every maximal ideal \(\mathfrak{m}\) of A;
\item
  one has dp \(_{A}\) (A/m) \textless{} n for every maximal ideal m of A.
\end{enumerate}

In particular, we have \(\mathrm{dh}(\Lambda)=\sup_{\mathfrak{m}}\mathrm{dp}_{\mathrm{A}}(\Lambda/\mathfrak{m})\) , where m ranges over the set of maximal ideals of A.

PROPOSITION 6. Let A be a Noetherian ring, N an A-module, and n an integer ≥ 0. The following conditions are equivalent:

\begin{enumerate}
\def\labelenumi{(\roman{enumi})}
\item
  we have \(\mathrm{di}_{\mathrm{A}}(\mathrm{N}) < n\) ;
\item
  for every prime ideal \(\mathfrak{p}\) of \(A\) , one has \(\mathrm{Ext}_{\mathrm{A}}^{n}(A/\mathfrak{p},\mathrm{N})=0\) ;
\item
  for every prime ideal \(\mathfrak{p}\) of A, we have \(\operatorname{Ext}_{\mathrm{A}_{\mathfrak{p}}}^{n}(\kappa(\mathfrak{p}),\mathrm{N}_{\mathfrak{p}})=0\) .
\end{enumerate}

If moreover the A-module N is finitely generated, these conditions are equivalent to:

\begin{enumerate}
\def\labelenumi{(\roman{enumi})}
\setcounter{enumi}{3}
\tightlist
\item
  for every maximal ideal m of A, we have \(\operatorname{Ext}_{\mathrm{A}}^{i}(\mathrm{A}/\mathfrak{m},\mathrm{N})=0\) for \(n\leqslant i\leqslant n+\operatorname{ht}(\mathfrak{m})\) .
\end{enumerate}

Observe that condition (iii) is equivalent to

(iii') for every prime ideal \(\mathfrak{p}\) of A, we have \(\operatorname{Ext}_{\mathrm{A}}^{n}(\mathrm{A}/\mathfrak{p},\mathrm{N})\otimes_{\mathrm{A}}\kappa(\mathfrak{p})=0\) .

Indeed, as \(\operatorname{Ext}_{\mathrm{A}}^{n}(\mathrm{A}/\mathfrak{p},\mathrm{N})\) is annihilated by \(\mathfrak{p}\) , the A-module \(\operatorname{Ext}_{\mathrm{A}}^{n}(\mathrm{A}/\mathfrak{p},\mathrm{N})\otimes_{\mathrm{A}}\kappa(\mathfrak{p})\) is isomorphic to \(\operatorname{Ext}_{\mathrm{A}}^{n}(\mathrm{A}/\mathfrak{p},\mathrm{N})\otimes_{\mathrm{A}}\mathrm{A}_{\mathfrak{p}}\) , hence to \(\operatorname{Ext}_{\mathrm{A}_{\mathfrak{p}}}^{n}(\kappa(\mathfrak{p}),\mathrm{N}_{\mathfrak{p}})\) ( \(n^{\circ}\) 2, prop. 2).

The implications (i)⇒(ii) and (i)⇒(iv) follow from prop. 1 of no. 1, and the implication (ii)⇒(iii') is clear.

\begin{enumerate}
\def\labelenumi{(\roman{enumi})}
\setcounter{enumi}{2}
\tightlist
\item
  (iii): for every A-module M, set T(M) = Ext \(_{A}^{n}\) (M,N). Suppose that (i) does not hold. Then (no. 1, prop. 1) there exists an ideal a of A such that T(A/a) ≠ 0. Let p be an ideal of A maximal among the ideals having this property; let us prove that p is prime. By IV, § 1, no. 4, th. 1 and th. 2, there exists a composition series (M \(_{i}\) ) \(_{0\leqslant i\leqslant m}\) of A/p such that each quotient M \(_{i}\) /M \(_{i+1}\) is isomorphic to a module A/p \(_{i}\) , where p \(_{i}\) (0 ≤ i ≤ m - 1) is a prime ideal containing p. If T(A/p \(_{i}\) ) were zero for all i, one would deduce by induction on i from the exact sequences
\end{enumerate}

\[
\mathrm{T} \left(\mathrm{M} _ {0} / \mathrm{M} _ {i}\right) \longrightarrow \mathrm{T} \left(\mathrm{M} _ {0} / \mathrm{M} _ {i + 1}\right) \longrightarrow \mathrm{T} \left(\mathrm{M} _ {i} / \mathrm{M} _ {i + 1}\right)
\]

that \(\mathrm{T}(\mathrm{M}_{0}/\mathrm{M}_{m})=\mathrm{T}(\mathrm{A}/\mathfrak{p})\) is zero, which is not the case. There therefore exists an index i such that \(\mathrm{T}(\mathrm{A}/\mathfrak{p}_{i})\neq0\) . By the maximality of p, we have \(p=p_{i}\) , so that p is prime.

Let \(x\) be an element of \(\mathrm{A} - \mathfrak{p}\) ; we deduce from the exact sequence

\[
0 \rightarrow \mathrm{A} / \mathfrak {p} \stackrel {{x _ {\mathrm{A} / \mathfrak {p}}}} {{\longrightarrow}} \mathrm{A} / \mathfrak {p} \longrightarrow \mathrm{A} / (\mathfrak {p} + x \mathrm{A}) \rightarrow 0
\]

an exact sequence

\[
\mathrm{T} (\mathrm{A} / (\mathfrak {p} + x \mathrm{A})) \longrightarrow \mathrm{T} (\mathrm{A} / \mathfrak {p}) \stackrel {{u}} {{\longrightarrow}} \mathrm{T} (\mathrm{A} / \mathfrak {p}),
\]

where \(u = \operatorname{Ext}_{\mathrm{A}}^{n}(x_{\mathrm{A}}, 1_{\mathrm{N}}) - x_{\mathrm{T(A/p)}}\) (A, X, p.~89, prop. 6). By the maximality of \(\mathfrak{p}\) , one has \(\mathrm{T(A / (\mathfrak{p} + xA))} = 0\) , so that the homomorphism \(u\) is injective. The nonzero A/p-module \(\mathrm{T(A/p)}\) is therefore torsion-free. This implies that \(\mathrm{T(A/p)} \otimes_{\mathrm{A}} \kappa(\mathfrak{p})\) is nonzero (A, II, p.~117, cor. 1), which contradicts (iii').

Suppose the A-module N is finitely generated, and let us prove that (iv) implies (iii). Let p be a prime ideal of A and m a maximal ideal of A containing p. There exists a saturated chain of prime ideals of A with endpoints p and m. The length r of this chain is less than ht(m); under hypothesis (iv), we therefore have

\[
\operatorname{Ext} _ {\mathrm{A} _ {\mathfrak {m}}} ^ {n + r} (\kappa (\mathfrak {m}), \mathrm{N} _ {\mathfrak {m}}) = \operatorname{Ext} _ {\mathrm{A}} ^ {n + r} (\mathrm{A} / \mathfrak {m}, \mathrm{N}) \otimes_ {A} \mathrm{A} _ {\mathfrak {m}} = 0.
\]

It then follows from Lemma 3 of § 1, No. 7 that we have \(\mathrm{Ext}_{\mathrm{A}_{\mathfrak{p}}}^{n}(\kappa(\mathfrak{p}),\mathrm{N}_{\mathfrak{p}})=0\), which proves (iii).

Remark 5.— Let N be a finitely generated A-module; the condition \(\operatorname{Ext}_{\mathrm{A}}^{n}(A/\mathfrak{m},\mathrm{N})=0\) for every maximal ideal \(\mathfrak{m}\) of A does not necessarily imply \(\operatorname{di}_{\mathrm{A}}(\mathrm{N})<n\). If, for example, the ring A is local and is not a Gorenstein ring (no. 7, def. 1), we have \(\operatorname{Ext}_{A}^{n}(A/\mathfrak{m},\mathrm{A})=0\) for \(n<\operatorname{prof}(\mathrm{A})\) but \(\operatorname{di}_{\mathrm{A}}(\mathrm{A})=+\infty\) .

